% Stage 8 restored fragment — 2.3.4 Retardation and causal propagation (archive lines 4420-4661)
\subsection{Retardation and causal propagation}
\label{sec:ch233}
Subsection~\ref{sec:ch233} supplies the definitive account of \emph{retarded} time-domain
Green operators, their causal support, and the kernel-level upgrade of the resolvent
equivalence Eq.~\eqref{eq:ch2.resolvent-green-equivalence} and harmonic theorem
Eq.~\eqref{eq:ch2.harmonic-equivalence-theorem} from Subsection~\ref{sec:ch223}
\cite{stratton1941}. Outgoing causality and Laplace--Fourier transforms are presupposed
through Eqs.~\eqref{eq:ch2.causal-operator-support}--\eqref{eq:ch2.fourier-phasor-identification};
the display below constructs the Green kernels that realize those statements.
Harmonic formulas Eqs.~\eqref{eq:ch2.free-space-scalar-green} and~\eqref{eq:ch2.free-space-dyadic-green}
are cited; they appear as $s=-\jj\omega$ evaluations of causal Laplace
transforms under Sommerfeld outgoing selection at infinity \cite{stratton1941}.

Let $\mathbf{j}_{\mathrm{imp}}(\mathbf{r},t)$ be causal with
\(\operatorname{supp}_{(\mathbf{r},t)}\subseteq\Omega_s\times[0,\infty)\),
\(\Omega_s\Subset\Omega\), and let \(\mathcal{B}_\Gamma\), \(\varepsilon\), and
\(\mu\) be time-independent. Initial data are homogeneous,
\(\mathbf{U}_0=\mathbf{0}\), unless noted for switch-on illustrations. Vacuum
regions use \(c_0=1/\sqrt{\mu_0\varepsilon_0}\) and \(k_0=\omega/c_0\) from
Eq.~\eqref{eq:ch2.helmholtz-symbol} \cite{stratton1941}. The equivalence class
excludes advanced kernels satisfying Eq.~\eqref{eq:ch2.advanced-kernel-violation}.
The causal scalar Green function for the wave operator
\((\nabla^{2}-\varepsilon_0\mu_0\,\partial_t^{2})\) in vacuum satisfies the
inhomogeneous wave equation below, with the time argument \(t-t'\) measuring elapsed
interval between source event and field observation (seconds, SI).
\begin{equation}
\label{eq:ch2.retarded-scalar-wave-operator}
\Big(\nabla^{2}-\varepsilon_0\mu_0\,\partial_t^{2}\Big)\,
G(\mathbf{r},\mathbf{r}',t-t')
=
-\delta(\mathbf{r}-\mathbf{r}')\,\delta(t-t'),
\end{equation}
with retarded causality $G=\mathbf{0}$ for $t-t'<0$: no signal precedes the light cone
\cite{stratton1941}. Time-harmonic kernels inherit this support; acausal branches appear
as standing-wave contamination in open domains \cite{harrington2001}.
The free-space retarded kernel concentrates its support on the light cone
\(t-t'=R/c_0\), where \(c_0=1/\sqrt{\mu_0\varepsilon_0}\) is the vacuum speed of light
(m\,s\(^{-1}\)) and \(R=|\mathbf{r}-\mathbf{r}'|\) is the geometric separation in metres.
The free-space outgoing solution is
\begin{equation}
\label{eq:ch2.retarded-scalar-green-time}
G_0(\mathbf{r},\mathbf{r}',t-t')
=
\frac{\delta\!\big(t-t'-R/c_0\big)}{4\pi R},
\qquad
R=|\mathbf{r}-\mathbf{r}'|,
\end{equation}
which transports impulse support on the light cone \(t-t'=R/c_0\)
\cite{stratton1941}. Equation~\eqref{eq:ch2.retarded-scalar-green-time} is the
time-domain counterpart of Eq.~\eqref{eq:ch2.free-space-scalar-green}: the locked
phasor field \(\exp(-\jj\omega t)\) selects the \(s=-\jj\omega\) Laplace transform
of Eq.~\eqref{eq:ch2.retarded-scalar-green-time} on the outgoing branch
\cite{harrington2001}. Physically, Eq.~\eqref{eq:ch2.retarded-scalar-green-time}
states that a unit impulse at \((\mathbf{r}',t')\) influences \(\mathbf{r}\) only
after travel time \(R/c_0\), enforcing finite propagation speed and
Eq.~\eqref{eq:ch2.causal-support-constraint} on potentials generated by scalar
convolution \cite{stratton1941}.
The retarded electric dyadic \(\overline{\mathbf{G}}_{E}(\mathbf{r},\mathbf{r}',t-t')\)
upgrades Eq.~\eqref{eq:ch2.electric-dyadic-convolution} to the time domain.
Time differentiation of the causal kernel enforces that the electric field responds only
after the impressed current has acted for a retarded interval; the convolution below is
the dyadic analogue of the scalar retarded integral.
The causal Maxwell field is
\begin{equation}
\label{eq:ch2.retarded-dyadic-convolution}
\E(\mathbf{r},t)
=
\E_{\mathrm{hom}}(\mathbf{r},t)
-\mu_0\,\partial_t\int_{0}^{t}\!\!\int_{\Omega}
\overline{\mathbf{G}}_{E}(\mathbf{r},\mathbf{r}',t-t')\cdot
\mathbf{j}_{\mathrm{imp}}(\mathbf{r}',t')\,d^{3}r'\,dt',
\end{equation}
which specializes Eq.~\eqref{eq:ch2.retarded-convolution-form} once the dyadic kernel
is identified \cite{stratton1941}. Retarded causality requires
\begin{equation}
\label{eq:ch2.retarded-dyadic-support}
\overline{\mathbf{G}}_{E}(\mathbf{r},\mathbf{r}',t-t')
=
\mathbf{0}\quad\text{for}\ t-t'<0,
\end{equation}
analogous to Eq.~\eqref{eq:ch2.retarded-kernel-support} \cite{stratton1941}. For
\(\mathbf{r}\notin\Omega_s\), Eq.~\eqref{eq:ch2.retarded-dyadic-convolution} produces
\(\E(\mathbf{r},t)=\mathbf{0}\) until the earliest retarded arrival time
\(t_{\mathrm{ret}}(\mathbf{r})\), defined as the infimum over all source events
\((\mathbf{r}',t')\) in the support of \(\mathbf{j}_{\mathrm{imp}}\) of the sum
\(t'+R/c_0\):
\begin{equation}
\label{eq:ch2.light-cone-arrival-time}
t_{\mathrm{ret}}(\mathbf{r})
=
\inf_{(\mathbf{r}',t')\in\operatorname{supp}\mathbf{j}_{\mathrm{imp}}}
\big\{t'\,+\,R/c_0\big\},
\end{equation}
which is the geometric light-cone support condition underlying
Eq.~\eqref{eq:ch2.causal-support-constraint} \cite{jackson1999}. The
harmonic dyad Eq.~\eqref{eq:ch2.free-space-dyadic-green} is not repeated here; it is
the \(s=-\jj\omega\) symbol of \(\overline{\mathbf{G}}_{E,0}\) under
Eqs.~\eqref{eq:ch2.equivalence-assumptions} and outgoing radiation condition
\cite{harrington2001}.
\begin{figure}[t]
\centering
\ChOneFigBlock{%
\begin{tikzpicture}[ch1fig, node distance=5mm and 7mm]
 \node[ch1box] (src) {$(\mathbf{r}',t')$\\source event};
 \node[ch1box, right=16mm of src] (cone) {Light cone\\$t-t'=R/c_0$};
 \node[ch1box, right=16mm of cone] (obs) {$(\mathbf{r},t)$\\observation};
 \node[ch1box, below=5mm of cone] (adv) {Advanced\\excluded};
 \draw[ch1flow] (src) -- (cone) -- (obs);
 \draw[ch1flow,dashed] (src) to[bend right=20] (adv);
\end{tikzpicture}%
}
\caption{Retarded support Eq.~\eqref{eq:ch2.retarded-dyadic-support}: field at
\((\mathbf{r},t)\) depends only on sources with \(t'\le t-R/c_0\); advanced
coupling is excluded.}
\label{fig:ch2.retarded-light-cone}
\end{figure}
~\eqref{eq:ch2.retarded-dyadic-convolution}.
Subsection~\ref{sec:ch232} treated coincidence in \(\mathbf{r}\); the display below
completes the causal picture by fixing the temporal ordering enforced by retarded support
\cite{stratton1941}. Regularized self-terms from
Eq.~\eqref{eq:ch2.mom-self-term-regularized} apply to harmonic convolutions; time
stepping discretizes Eq.~\eqref{eq:ch2.retarded-dyadic-convolution} with the same
kernel class \cite{harrington2001}. Let \(\widetilde{\overline{\mathbf{G}}}_{E}(\mathbf{r},\mathbf{r}',s)\) denote the
Laplace transform of \(\overline{\mathbf{G}}_{E}(\mathbf{r},\mathbf{r}',t)\) in
\(t-t'\) with \(\Re\{s\}\) large enough for convergence.
The complex variable \(s\) (s\(^{-1}\)) encodes temporal frequency; evaluation at
\(s=-\jj\omega\) recovers the harmonic dyad used throughout the phasor convention.
Causal support
Eq.~\eqref{eq:ch2.retarded-dyadic-support} implies
\(\widetilde{\overline{\mathbf{G}}}_{E}\) is analytic in the right half-plane and,
under outgoing radiation condition,
\begin{equation}
\label{eq:ch2.retarded-laplace-harmonic-link}
\overline{\mathbf{G}}_{E,\omega}(\mathbf{r},\mathbf{r}')
=
\eval{\widetilde{\overline{\mathbf{G}}}_{E}(\mathbf{r},\mathbf{r}',s)}_{s=-\jj\omega},
\end{equation}
which is the kernel form of Eq.~\eqref{eq:ch2.resolvent-green-equivalence}
\cite{stratton1941}. Applying the Laplace transform to
Eq.~\eqref{eq:ch2.retarded-dyadic-convolution} with causal
\(\mathbf{j}_{\mathrm{imp}}\) yields
\begin{equation}
\label{eq:ch2.laplace-dyadic-convolution}
\widetilde{\E}(\mathbf{r},s)
=
\widetilde{\E}_{\mathrm{hom}}(\mathbf{r},s)
-\mu_0 s\int_{\Omega}
\widetilde{\overline{\mathbf{G}}}_{E}(\mathbf{r},\mathbf{r}',s)\cdot
\widetilde{\mathbf{j}}(\mathbf{r}',s)\,d^{3}r',
\end{equation}
with \(\widetilde{\mathbf{j}}\) from Eq.~\eqref{eq:ch2.laplace-source-transform}
\cite{stratton1941}. On \(s=-\jj\omega\) under Eqs.~\eqref{eq:ch2.equivalence-assumptions},
Eq.~\eqref{eq:ch2.laplace-dyadic-convolution} reduces to
Eq.~\eqref{eq:ch2.electric-dyadic-convolution} with
\(\overline{\mathbf{G}}_{E,\omega}\) from Eq.~\eqref{eq:ch2.retarded-laplace-harmonic-link}
and \(\Jimp=\eval{\widetilde{\mathbf{j}}(\mathbf{r},s)}_{s=-\jj\omega}\)
\cite{harrington2001}. The Chapter~\ref{ch:02} causal propagation theorem at kernel
level is
\begin{equation}
\label{eq:ch2.retarded-harmonic-kernel-theorem}
\mathcal{S}_\omega[\Jimp]
=
\eval{\mathcal{S}_t[\mathbf{j}_{\mathrm{imp}}](\mathbf{r},t)}_{\text{harmonic at }\omega}
\quad\Longleftrightarrow\quad
\overline{\mathbf{G}}_{E,\omega}
=
\eval{\widetilde{\overline{\mathbf{G}}}_{E}}_{s=-\jj\omega}
\end{equation}
for the same \(\mathcal{B}_\Gamma\) and outgoing branch, with \(\mathcal{S}_t\) the
time-domain map Eq.~\eqref{eq:ch2.time-domain-solution-map} \cite{stratton1941}.
Equation~\eqref{eq:ch2.retarded-harmonic-kernel-theorem} does not replace
Eq.~\eqref{eq:ch2.harmonic-equivalence-theorem}; it identifies the Green operators
whose resolvent relation Eq.~\eqref{eq:ch2.resolvent-green-equivalence} implements
\cite{harrington2001}.
Advanced kernels and branch failure. Advanced dyads satisfying Eq.~\eqref{eq:ch2.advanced-kernel-violation} produce
Laplace symbols that cannot equal
Eq.~\eqref{eq:ch2.retarded-laplace-harmonic-link} for any causal
\(\mathbf{j}_{\mathrm{imp}}\) \cite{stratton1941}. Time-harmonic solvers that
implicitly use incoming Green functions therefore fail
Eq.~\eqref{eq:ch2.retarded-harmonic-kernel-theorem} even when interior Maxwell rows
vanish \cite{harrington2001}. Exterior problems require outgoing radiation condition consistent
with Section~\ref{sec:ch24}; interior cavity problems may require loss or quotienting
as in Eq.~\eqref{eq:ch2.coercivity-quotient} before a causal Laplace contour can be
closed \cite{stratton1941}. For finite-bandwidth excitation, the harmonic slice at fixed \(\omega\) is insufficient.
Causal fields superpose as
\begin{equation}
\label{eq:ch2.broadband-retarded-synthesis}
\E(\mathbf{r},t)
=
\int_{-\infty}^{\infty}
\widetilde{\E}(\mathbf{r},-\jj\omega)\,e^{-\jj\omega t}\,\frac{d\omega}{2\pi},
\end{equation}
with \(\widetilde{\E}(\mathbf{r},-\jj\omega)\) the Fourier transform induced by
Eq.~\eqref{eq:ch2.laplace-dyadic-convolution} on causal data
\cite{stratton1941}. Equation~\eqref{eq:ch2.broadband-retarded-synthesis}
is the inverse of the single-frequency reduction used in the phasor convention; dispersive
\(\varepsilon(\omega)\), \(\mu(\omega)\) modify
\(\widetilde{\overline{\mathbf{G}}}_{E}(\mathbf{r},\mathbf{r}',s)\) and are deferred
beyond the locked single-\(\omega\) harmonic band except where loss enters as
\(\Im\{k_0\}\neq 0\) \cite{harrington2001}. Within the harmonic band,
Eq.~\eqref{eq:ch2.steady-state-harmonic-limit} describes the long-time limit after
switch-on when outgoing radiation or damping removes transients
\cite{stratton1941}.
Because \(\Pi_{\mathrm{rad}}\) commutes with harmonic extraction under the ordering
assumed in Subsection~\ref{sec:ch223},
\begin{equation}
\label{eq:ch2.retarded-radiation-kernel-composition}
\mathcal{R}_t[\mathbf{j}_{\mathrm{imp}}]
=
\Pi_{\mathrm{rad}}\circ\mathcal{S}_t[\mathbf{j}_{\mathrm{imp}}],
\qquad
\mathcal{R}_\omega[\Jimp]
=
\Pi_{\mathrm{rad}}\circ\mathcal{S}_\omega[\Jimp],
\end{equation}
with \(\mathcal{R}_\omega\) from Eq.~\eqref{eq:ch2.radiation-time-frequency-equiv}
\cite{stratton1941}. Causality does not promote evanescent content to
radiative export: Eq.~\eqref{eq:ch2.regime-operator-screen} remains the transport criterion
after time-domain evolution \cite{stratton1941}. Retarded propagation ensures support
constraints only; spectral and regime qualifications from Subsection~\ref{sec:ch224} are
additional \cite{harrington2001}.
Equations~\eqref{eq:ch2.retarded-dyadic-convolution}--\eqref{eq:ch2.retarded-harmonic-kernel-theorem}
require Eqs.~\eqref{eq:ch2.equivalence-assumptions}, outgoing or passive completion on
exterior domains, and \(\mathbf{j}_{\mathrm{imp}}\) in the causal class of
Subsection~\ref{sec:ch221}. They fail for time-varying \(\mathcal{B}_\Gamma(t)\),
moving media, and nonlinear sources. Coincidence and self-term evaluation of
\(\overline{\mathbf{G}}_{E,\omega}\) remains governed by Subsection~\ref{sec:ch232}
even when Eq.~\eqref{eq:ch2.retarded-dyadic-convolution} is discretized in time
\cite{harrington2001}. Retarded support is therefore the temporal counterpart of the spatial
singular split in Subsection~\ref{sec:ch232}: causality fixes when propagation may act, while
regularization fixes how strongly a source couples to itself at coincidence
\cite{stratton1941}. Harmonic kernels $G$ and $\overline{\mathbf{G}}_{E}$ are therefore
retarded response functions in space and time; constructive inversion on composite domains asks
which resolvent branch those responses realize when interfaces and exterior completions are
assembled \cite{stratton1941}.
